% This document was generated by an LLM.

\documentclass{essay}

\title{Asymptotic Analysis: Big-$O$ and Little-$o$}
\subtitle{Landau notation, asymptotic equivalence, and their algebra, for calculus and real analysis}
\author{}
\date{}

\begin{document}
\maketitle

\section{Setup}

Throughout, $f$ and $g$ are real-valued functions defined near a fixed base point $a$, where $a$ may be a finite real number or $\pm\infty$. ``Near $a$'' means on some deleted neighborhood of $a$ if $a\in\R$, or on some ray $(M,\infty)$ or $(-\infty,M)$ if $a=\pm\infty$. All statements below are local: a function can be $O(g)$ as $x\to 0$ and something completely different as $x\to\infty$, so the base point is always part of the claim, even when left implicit.

\begin{note}
When a limit definition below divides by $g(x)$, assume $g(x)\neq 0$ near $a$ (except possibly at $a$ itself).
\end{note}

\section{The Landau symbols}

\begin{definition}[Little-$o$]
$f(x) = o\big(g(x)\big)$ as $x\to a$ means: for every $\varepsilon>0$ there is a neighborhood of $a$ on which
\[
  |f(x)| \le \varepsilon\,|g(x)|.
\]
Equivalently, if $g$ is eventually nonzero near $a$,
\[
  \lim_{x\to a} \frac{f(x)}{g(x)} = 0.
\]
Read ``$f$ is of strictly smaller order than $g$ near $a$,'' or ``$f$ is negligible compared to $g$.''
\end{definition}

\begin{definition}[Big-$O$]
$f(x) = O\big(g(x)\big)$ as $x\to a$ means: there exist $C>0$ and a neighborhood of $a$ on which
\[
  |f(x)| \le C\,|g(x)|.
\]
Equivalently, if $g$ is eventually nonzero near $a$, the ratio $f(x)/g(x)$ is \emph{bounded} near $a$ --- it need not converge.
Read ``$f$ is of order at most $g$ near $a$.''
\end{definition}

\begin{definition}[Big-$\Theta$]
$f(x) = \Theta\big(g(x)\big)$ as $x\to a$ means $f=O(g)$ and $g=O(f)$: there exist $c,C>0$ with
\[
  c\,|g(x)| \le |f(x)| \le C\,|g(x)|
\]
near $a$. Read ``$f$ and $g$ have the same order'' --- a symmetric relation, unlike $O$ and $o$.
\end{definition}

\begin{definition}[Asymptotic equivalence]
$f(x) \sim g(x)$ as $x\to a$ means
\[
  \lim_{x\to a} \frac{f(x)}{g(x)} = 1.
\]
This is strictly stronger than $\Theta$: it fixes the constant of proportionality to be exactly $1$ in the limit, not merely bounded above and below.
\end{definition}

\begin{remark}
These are statements about \emph{one function belonging to a set of functions}: $O(g)$ denotes the set of all $f$ satisfying the bound above, and ``$f=O(g)$'' is shorthand for ``$f\in O(g)$.'' The equals sign is an abuse of notation inherited from long-standing convention, not a genuine equality --- it is not symmetric ($x=O(x^2)$ as $x\to 0$, but writing $O(x^2)=x$ is meaningless), and two different functions can both equal $O(g)$ without equaling each other.
\end{remark}

\section{Basic algebra}

Fix the base point $a$. These hold for functions compared near $a$.

\begin{itemize}
  \item \textbf{$o$ implies $O$.} If $f=o(g)$ then $f=O(g)$ (take $C=1$ and shrink the neighborhood). The converse fails: $f=O(g)$ does not imply $f=o(g)$.
  \item \textbf{Transitivity.} $f=O(g)$ and $g=O(h)$ imply $f=O(h)$; likewise for $o$, and for mixed chains ($f=o(g)$, $g=O(h)$ $\Rightarrow$ $f=o(h)$).
  \item \textbf{Constant multiples are absorbed.} $O(cg)=O(g)$ and $o(cg)=o(g)$ for any constant $c\neq 0$. In particular $c\cdot O(g) = O(g)$.
  \item \textbf{Sum rule (absorption).} $O(g)+O(g)=O(g)$, and more generally $O(g)+O(h) = O(\max(|g|,|h|))$. The same holds for $o$. A sum is governed by its dominant term: $o(g)+O(g)=O(g)$.
  \item \textbf{Product rule.} $O(f)\cdot O(g) = O(fg)$, and $o(f)\cdot O(g) = o(fg)$. Multiplying two little-$o$ terms gives little-$o$ of the product: $o(f)\cdot o(g) = o(fg)$.
  \item \textbf{Powers.} If $f=O(g)$ with $f,g>0$ near $a$, then $f^k = O(g^k)$ for $k>0$; similarly for $o$.
  \item \textbf{No cancellation.} $O(g)-O(g)$ is \emph{not} $0$; it is again $O(g)$. Landau symbols name an upper bound on size, not a specific function, so ``subtracting'' two unspecified bounded quantities does not cancel them. Never manipulate an $O$- or $o$-expression as if it were an ordinary algebraic term standing for one fixed function.
  \item \textbf{Composition is not automatic.} $f=O(g)$ as $x\to a$ does \emph{not} generally give $h(f(x)) = O(h(g(x)))$ for an arbitrary $h$; this needs extra hypotheses on $h$ (e.g.\ $h$ Lipschitz, or monotone with controlled growth).
\end{itemize}

\section{Taylor's theorem in Landau notation}

Landau notation is the natural language for truncating a Taylor expansion: it records how fast the error shrinks without writing the error term explicitly.

\begin{theorem}[Peano remainder]
If $f$ is $n$ times differentiable at $a$, then as $x\to a$,
\[
  f(x) = \sum_{k=0}^{n} \frac{f^{(k)}(a)}{k!}(x-a)^k + o\big((x-a)^n\big).
\]
\end{theorem}

\begin{theorem}[Lagrange / big-$O$ remainder]
If $f$ is $(n+1)$ times differentiable on an interval containing $a$ and $x$, with the $(n+1)$-st derivative bounded there, then
\[
  f(x) = \sum_{k=0}^{n} \frac{f^{(k)}(a)}{k!}(x-a)^k + O\big((x-a)^{n+1}\big)
\]
as $x\to a$ --- one order sharper than the Peano form, at the cost of the stronger smoothness hypothesis.
\end{theorem}

\begin{example}[Standard expansions as $x\to 0$]
\begin{align*}
  e^x &= 1 + x + \frac{x^2}{2} + O(x^3) \\
  \sin x &= x - \frac{x^3}{6} + o(x^3) \\
  \cos x &= 1 - \frac{x^2}{2} + O(x^4) \\
  \ln(1+x) &= x - \frac{x^2}{2} + O(x^3) \\
  (1+x)^r &= 1 + rx + \binom{r}{2}x^2 + o(x^2) \qquad (r\in\R)
\end{align*}
\end{example}

\begin{remark}
The exponent inside the Landau term is the order of the \emph{first omitted} term, not the number of terms kept. Writing $o(x^3)$ after a cubic term and $O(x^3)$ after a quadratic term are both legitimate but say different things about where the next correction enters.
\end{remark}

\section{Asymptotic equivalence: the finer tool}

\begin{itemize}
  \item $f\sim g$ as $x\to a$ \emph{iff} $f - g = o(g)$, i.e.\ $f = g + o(g) = g\big(1+o(1)\big)$.
  \item \textbf{Products and quotients behave well.} If $f_1\sim g_1$ and $f_2\sim g_2$, then $f_1f_2 \sim g_1g_2$ and $f_1/f_2 \sim g_1/g_2$. Powers follow too: $f^r \sim g^r$ for fixed $r$ (where defined).
  \item \textbf{Sums do not.} $f_1\sim g_1$ and $f_2\sim g_2$ does \emph{not} imply $f_1+f_2 \sim g_1+g_2$ --- the leading terms can cancel. Example, as $x\to\infty$: take $g_1(x)=x$, $f_1(x)=x+1$ (so $f_1\sim g_1$), and $g_2(x)=-x$, $f_2(x)=-x+\sqrt{x}$ (so $f_2\sim g_2$, since $f_2/g_2 = 1-1/\sqrt{x}\to 1$). Then $g_1+g_2\equiv 0$, while $f_1+f_2 = 1+\sqrt{x}\to\infty$ --- so $f_1+f_2$ cannot be asymptotically equivalent to $g_1+g_2$. \emph{Asymptotic equivalence is safe to use inside products and quotients, but never inside a sum or difference without checking that the leading terms survive cancellation.}
  \item \textbf{Equivalence controls the limit of ratios, not the size of the difference.} $f\sim g$ says nothing about $f-g\to 0$; it only says $f-g$ is small \emph{relative to $g$}. E.g.\ $x^2+x \sim x^2$ as $x\to\infty$, while the difference $x$ itself diverges.
\end{itemize}

\begin{example}[Stirling's formula]
\[
  n! \sim \sqrt{2\pi n}\,\left(\frac{n}{e}\right)^n \qquad (n\to\infty).
\]
The ratio of the two sides tends to $1$, even though their difference grows without bound --- the hallmark of $\sim$ as distinct from a literal approximation.
\end{example}

\section{Comparing growth at infinity}

A standard hierarchy of orders of infinity, each term $o$ of the next as $x\to\infty$ (for constants $p>0$, $b>1$, and $0<c<1$):
\[
  1 \;\prec\; \ln x \;\prec\; x^{\varepsilon} \;\prec\; x^{p} \;\prec\; x^{p}\ln x \;\prec\; x^{p+\varepsilon} \;\prec\; e^{x^{c}} \;\prec\; b^{x} \;\prec\; x! \;\prec\; x^{x},
\]
where $f\prec g$ abbreviates $f=o(g)$, $\varepsilon>0$ is arbitrarily small, and $b_1^x \prec b_2^x$ whenever $1<b_1<b_2$. This hierarchy is the usual first resort for limit-comparison arguments (e.g.\ deciding convergence of $\int^{\infty} f$ or $\sum f(n)$ by comparing $f$ against a power of $x$ or an exponential).

\begin{example}
$\displaystyle\lim_{x\to\infty}\frac{\ln x}{x^{0.01}} = 0$ follows immediately from $\ln x \prec x^{\varepsilon}$ above, without L'Hôpital.
\end{example}

\section{Common pitfalls}

\begin{itemize}
  \item \textbf{Wrong direction of comparison.} $f=O(g)$ bounds $f$ \emph{above} in size; it says nothing about $g=O(f)$. Confusing ``at most the order of'' with ``the same order as'' is the single most common error.
  \item \textbf{$O$ does not imply the limit of the ratio exists.} $\sin x = O(1)$ as $x\to\infty$, but $\sin x$ has no limit and $\sin x \not\sim c$ for any constant $c$.
  \item \textbf{The base point matters.} $x^2 = O(x)$ fails as $x\to\infty$ but holds as $x\to 0$; always state (or infer from context) which limit is in force before writing a Landau expression.
  \item \textbf{$o(1)$ and $O(1)$ are shorthand for ``tends to $0$'' and ``stays bounded.''} These are the two most frequent instances and worth recognizing on sight: $f=o(1) \iff f(x)\to 0$; $f=O(1) \iff f$ is bounded near $a$.
  \item \textbf{Landau notation is not a substitute for a rigorous bound when one is needed.} It is a bookkeeping device for tracking leading-order behavior; when a problem asks for an explicit constant or an explicit neighborhood, the $O$ has to be unpacked back into its $\varepsilon$--$\delta$ (or $C$--neighborhood) form.
\end{itemize}

\end{document}
